% 351_PDG_Modellabhaengigkeit_En_ch.tex
% Johann Pascher, 5 September 2026
\providecommand{\BM}{\textbf{[B]}}
\providecommand{\KM}{\textbf{[K]}}
\providecommand{\SM}{\textbf{[S]}}
\providecommand{\QM}{\textbf{[Q]}}

\phantomsection
\addcontentsline{toc}{chapter}{Model Dependence of PDG Comparison Values}

\section*{Status markers}
\begin{center}
\begin{tabular}{@{}lp{10.5cm}@{}}
\textbf{[K]} & \emph{Confirmed} --- derived from $\xi$, numerically verified. \\[4pt]
\textbf{[B]} & \emph{Proved} --- algebraically established. \\[4pt]
\textbf{[S]} & \emph{Speculative} --- open. \\[4pt]
\textbf{[Q]} & \emph{Source} --- corpus-external primary source, verified here. \\
\end{tabular}
\end{center}

\section*{Abstract}

FFGFT predictions are compared against PDG and CODATA values. This document examines how those comparison values actually arise: which raw observable underlies each, which theory (QED, Standard Model) enters the extraction, and where irrational numbers already appear in the extraction formulae themselves.

Main findings: (1)~$m_e$ and $m_\mu/m_e$ are the least model-dependent anchors, but $m_\mu/m_e$ is QED-dependent and derives from a 27-year-old muonium HFS measurement whose theory uncertainty CODATA underestimates by a factor of 5--10 (Eides 2026). (2)~$\alpha$ carries an unresolved $5\sigma$ tension between the two best recoil measurements. (3)~$v=246$~GeV is never measured; it is the SM definition $v\equiv(\sqrt2\,G_F)^{-1/2}$. (4)~$G_F$ carries a theory correction $\Delta q=-0.44\,\%$ with a $4.7\sigma$ hadronic tension. (5)~$m_\tau$ has two measurement methods differing by 0.18~MeV; FFGFT lies $0.4\sigma$ from the PDG average. (6)~The magnitude of all theory uncertainties ($10^{-7}$ to $10^{-10}$) does not explain the $\sim1\,\%$ residual in $(m_\mu/m_e)^2$ --- it remains physical. (7)~Errors accumulate along the derivation chain; every prediction requires its full provenance path.

All 21 primary sources were verified on 5~September 2026.

\section{The question}

The corpus compares FFGFT predictions with PDG values and finds residuals at the sub-percent to percent level. Are PDG values theory-neutral, or do they themselves depend on the Standard Model --- such that part of the residual arises on the extraction side rather than the FFGFT side?

The FFGFT foundational statement: only ratios free of irrational numbers are exact; SI absolute values are structurally approximate (Docs.~241, 085). If PDG values themselves contain irrational factors in their extraction, the same limitation applies to the comparison side.

\section{How empirical values arise}

No PDG or CODATA value is a raw measurement. Each passes through the same chain:
\begin{center}
raw observable $\to$ theory formula $\to$ extracted parameter $\to$ CODATA adjustment $\to$ PDG value
\end{center}

\subsection{The CODATA adjustment}

CODATA publishes recommended values every four years, produced by a joint least-squares adjustment over roughly 300 input data linked by theory relations (QED, Rydberg formula, kinematics). The values are therefore \emph{correlated}: a change in one input shifts all dependent values. Where theory uncertainties enter, they are partly treated as adjustment parameters ($\delta_\text{th}$ in the observational equations) rather than as errors --- this is the Eides finding for muonium HFS \QM.

\subsection{Electron mass $m_e$}

\textbf{Raw observable:} cyclotron frequency of a hydrogen-like ion ($^{12}$C$^{5+}$, $^{28}$Si$^{13+}$, $^{20}$Ne$^{9+}$) in a Penning trap, and the Larmor precession frequency of the bound electron. The ratio is a pure number.

\textbf{Theory formula:} bound $g$-factor from QED bound-state theory. From the frequency ratio and $g_\text{bound}$ follows $m_e/m_\text{ion}$; with the ion mass in u follows $m_e$ in u; conversion to MeV via u\,$\to$\,kg\,$\to$\,MeV.

\textbf{Theory dependence:} QED for $g_\text{bound}$, small and well controlled. Current: Heiße et al.\ (2023) at 0.1~ppb \QM.

\subsection{Muon-electron mass ratio $m_\mu/m_e$}

\textbf{Raw observable:} ground-state hyperfine splitting of muonium (a bound $\mu^+e^-$ state) --- a microwave resonance at 4.463~GHz --- and Zeeman splitting in a magnetic field.

\textbf{How $m_\mu/m_e$ is measured.} The muon is not placed on a scale.
From the measured HFS frequency $\Delta\nu$ and the QED formula for the Fermi frequency
\begin{equation}
\nu_F=\tfrac{16}{3}\,\alpha^2\,c\,R_\infty\,\frac{m_e}{m_\mu}\Big(\frac{m_r}{m_e}\Big)^3
\end{equation}
(corrections: QED to $O(\alpha^3)$, recoil in $m_e/m_\mu$, weak $Z$ exchange, hadronic VP) $m_\mu/m_e$ follows. Two consequences:
\begin{itemize}\itemsep2pt
\item $m_\mu/m_e$ \emph{is not theory-free}: the extracted value contains assumptions about
  $\alpha$, higher-order QED loop corrections, and hadronic vacuum polarisation.
  Since $\nu_F\propto\alpha^2$, $\alpha$ enters implicitly --- anchor and $\alpha$ are not independent.
\item \emph{CODATA overstates the real precision}: CODATA 2022
  gives $\sigma_\text{rel}(m_\mu/m_e)=2.2\times10^{-8}$. The actual uncertainty of the
  HFS extraction (experiment + QED theory from Eides 2026 + $\alpha$ tension)
  is $\approx1.2\times10^{-7}$ --- a factor of 5--6 larger. \QM{}
\end{itemize}

\textbf{Theory dependence:} high. Sole precision measurement: LAMPF 1999 (Liu et al.). Eides (2026) shows that CODATA lists the measured value 4\,463\,302\,776(51)~Hz as the theoretical value; the actual theory uncertainty is 271--515~Hz \QM.

\subsection{Fine-structure constant $\alpha$}

Two routes that disagree at $5\sigma$ \QM:
\begin{itemize}\itemsep2pt
\item \textbf{Atom recoil:} $\alpha^2=2R_\infty\,(h/m_\text{atom})\,(m_\text{atom}/m_e)/c$. Rb (Morel et al.\ 2020): $\alpha^{-1}=137.035\,999\,206(11)$. Cs (Parker et al.\ 2018): $137.035\,999\,046(27)$. Difference $1.6\times10^{-7}$, $>5\sigma$.
\item \textbf{Electron $g-2$:} $a_e$ at $10^{-13}$ (Fan et al.\ 2023) plus 5-loop QED (12\,672 diagrams, Aoyama et al.). Gives $\alpha^{-1}=137.035\,999\,166(15)$ --- between Rb and Cs.
\end{itemize}
The CODATA value 137.035\,999\,177(21) is an average covering none of the measurements. FFGFT: $\alpha^{-1}=3700/27=137.037\,037\ldots$ (Doc.~338, R101), deviation 7.6~ppm --- 100$\times$ larger than the Rb/Cs tension.

\subsection{Tau mass $m_\tau$ --- two methods}

\begin{center}
\begin{tabular}{@{}llll@{}}
\toprule
Method & Experiment & Value (MeV) & Theory assumption \\
\midrule
Threshold scan & BESIII 2014 & $1776.91\pm0.12\,{}^{+0.10}_{-0.13}$ & QED cross-section \\
Pseudomass & Belle II 2023 & $1777.09\pm0.08\pm0.11$ & $m_{\nu_\tau}=0$ \\
\textbf{PDG average 2024} & --- & $\mathbf{1776.93\pm0.09}$ & --- \\
\bottomrule
\end{tabular}
\end{center}

\textbf{Threshold scan:} $\sigma(e^+e^-\to\tau^+\tau^-)$ versus beam energy near $2m_\tau\approx3.55$~GeV; $m_\tau$ as fit parameter of the QED threshold curve. \textbf{Pseudomass:} at 10.6~GeV, the kinematic edge of the invariant mass in $\tau\to3\pi\nu$; assumes $m_{\nu_\tau}=0$.

\textbf{For FFGFT:} $m_\tau=1776.97$~MeV (pred) lies 0.04~MeV $=0.4\sigma$ from the PDG average. The methods differ by 0.18~MeV. The falsification test is undecided; Belle II with $>1$~ab$^{-1}$ should reach $\sigma<0.05$~MeV.

\subsection{Fermi constant $G_F$ and vacuum expectation value $v$}

\textbf{Raw observable for $G_F$:} muon decay time distribution (MuLan, $10^{12}$ decays), $\tau_\mu=2.196\,980\,3(22)$~µs \QM.

\textbf{Theory formula} (Eberhart et al.\ 2026 \QM):
\begin{equation}
\tau_\mu^{-1}=\frac{G_F^2\,m_\mu^5}{192\pi^3}\,(1+\Delta q),\qquad \Delta q=(-4\,384\,678\pm34)\times10^{-9}\approx-0.44\,\%.
\end{equation}
$\Delta q$ contains QED to $O(\alpha^3)$, hadronic vacuum polarisation, electron-mass effects. Purely electroweak corrections are absorbed into $G_F$ by definition (Sirlin 1980). The hadronic component carries a $4.7\sigma$ tension (CMD-3 vs.\ older data).

\textbf{Raw observable for $v$:} none. $v\equiv(\sqrt2\,G_F)^{-1/2}=246.22$~GeV follows from $M_W=gv/2$ and $G_F/\sqrt2=g^2/(8M_W^2)$. \textbf{$v$ exists only within the Standard Model.} Its numerical value is the muon lifetime in SM clothing.

\textbf{For FFGFT:} Doc.~006/R104 gives $v=248.3$~GeV (without $K_\text{frak}$) or 245.6~GeV (with $K_\text{frak}=74/75$). Comparison with 246.22~GeV is a comparison of \emph{FFGFT structure against SM definition}, not against an observable. Only $\tau_\mu$ and $m_\mu$ are measured.

\subsection{Gravitational constant $G$}

Torsion balance (Li et al.\ 2018, 11.6~ppm) or atom interferometry (Rosi et al.\ 2014, 150~ppm). Theory: Newton. The only model-free value, but the least precise: CODATA $6.67430(15)\times10^{-11}$, 22~ppm \QM.

\subsection{Neutrino mass-squared differences}

Event rates versus $L/E$, three-flavour fit (NuFIT 6.0 \QM). $\Delta m^2_{31}$ depends on whether SK atmospheric data are included: ratio 33.7 (with) vs.\ 33.9 (without). FFGFT $360/11=32.73$; measurement uncertainty 1--3\,\% $\approx$ FFGFT residual. Not discriminating.

\subsection{Overview}

\begin{center}
\begin{tabular}{@{}llll@{}}
\toprule
Quantity & Raw observable & Theory in extraction & Convention \\
\midrule
$m_e$ & frequency ratio & QED ($g_\text{bound}$) & --- \\
$m_\mu/m_e$ & HFS frequency & QED, weak, hadr. & --- \\
$\alpha$ (Rb/Cs) & interferometer phase & $R_\infty$ (QED) & --- \\
$\alpha$ ($g-2$) & frequency ratio & 5-loop QED & --- \\
$m_\tau$ & rates vs.\ $E$ / edge & QED threshold / $m_{\nu_\tau}=0$ & beam calibration \\
$G_F$ & decay times & Fermi + QED + hadr. & Sirlin absorption \\
$v$ & --- & --- & $v\equiv(\sqrt2G_F)^{-1/2}$ \\
$G$ & angle / acceleration & Newton & --- \\
$\Delta m^2$ & rates vs.\ $L/E$ & 3-flavour & with/without SK-atm \\
\bottomrule
\end{tabular}
\end{center}

No value is theory-free. For FFGFT this means: comparison with PDG values is comparison with ``QED + convention'', not with raw observations.

\section{Irrational factors in the extraction formulae}

The FFGFT foundational statement --- ratios of rationals are exact, SI absolute values carry the approximation of the projection $S^1_m\to\mathbb R_t$ (Docs.~306, 333) --- applies to the SM extraction formulae themselves:

\begin{center}
\begin{tabular}{@{}lll@{}}
\toprule
Quantity & Formula & Irrational factor \\
\midrule
$G_F$ & $\tau_\mu^{-1}=G_F^2m_\mu^5/(192\pi^3)$ & $\pi^3$ \\
$v$ & $v=(\sqrt2\,G_F)^{-1/2}$ & $\sqrt2$ \\
$\alpha$ (recoil) & $\alpha^2=2R_\infty h/(mc)$ & $R_\infty\propto\alpha^2$: implicitly circular \\
HFS & $\nu_F\propto\alpha^2$ & $\alpha^2$ \\
$\Delta q$ & expansion in $\alpha/\pi$ & $\pi$ in every term \\
\bottomrule
\end{tabular}
\end{center}

``$G_F$ measured'' means: $\tau_\mu$ measured, divided by $192\pi^3$, square root taken. ``$v$ measured'' additionally means: divided by $\sqrt2$, square root again. The irrational factors come from the SM formulation (phase-space integral, Higgs normalisation), not from measurement. The only convention-free comparison would be at the level of $\tau_\mu=2.196\,980\,3$~µs --- which FFGFT would have to predict directly \SM.

\section{The residual in $(m_\mu/m_e)^2$: exact computation}

FFGFT bare (Doc.~338): $(m_\mu/m_e)^2=43200$ exactly, an integer, $=|\text{GF}(9)^*|^2\cdot5^2\cdot|\text{GF}(27)|$. With CODATA 2022 in exact rational arithmetic:
\begin{align}
(m_\mu/m_e)^2_\text{PDG} &= 42753.12 \\
\text{residual} &= 446.88\quad(1.045\,\%) \\
\sigma_\text{rel}\big((m_\mu/m_e)^2\big) &= 4.35\times10^{-8}\quad\Rightarrow\quad \text{residual}=240\,000\,\sigma .
\end{align}

Possible error sources, computed:
\begin{itemize}\itemsep2pt
\item CODATA uncertainty: $\pm0.002$ --- factor $2\times10^5$ too small.
\item Eides correction to HFS theory uncertainty (515 instead of 51~Hz): $\pm0.010$ --- residual remains $45\,000\,\sigma$.
\item Choice of $\alpha$ in the HFS anchor (FFGFT $3700/27$ instead of CODATA): shifts $(m_\mu/m_e)^2$ by $-1.3$ --- factor 500 too small.
\item To reach 43200, $m_\mu$ would need to be 551~keV higher (measurement error 2.3~eV), or $m_e$ 2.65~keV lower (error 0.15~meV), or the HFS 23~MHz lower (error 51~Hz).
\end{itemize}

\textbf{Conclusion \textbf{[B]}:} No measurement error explains the residual. It is physical and belongs to the theory: fractal-recursive correction (Docs.~295, 146, 338) and projection approximation (Docs.~306, 333). In the square, without square root: $43200\cdot(1-77.6\,\xi)=42753$; $K_\text{frak}=1-100\xi$ over-corrects by $22.4\,\xi$. Which power of $K_\text{frak}$ acts on the quadratic identity is not fixed in the corpus \SM.


\subsection*{Rational or irrational: where the comparison is clean}

The discussion above compared three Galois identities against PDG values.
Not all comparisons are equally clean:

\begin{center}
\begin{tabular}{@{}llll@{}}
\toprule
Comparison & Galois value & Type & Residual \\
\midrule
$(m_\mu/m_e)^2$ vs.\ $43200$ & $43200$ (integer) & rational & $-1.034\,\%$ \\
$m_\mu/m_e$ vs.\ $\sqrt{43200}$ & $120\sqrt{3}$ & \textbf{irrational!} & $-0.519\,\%$ \\
$m_e\cdot m_\mu$ vs.\ $54~\text{MeV}^2$ & $54$ (integer) & rational & $-0.016\,\%$ \\
\bottomrule
\end{tabular}
\end{center}

The linear comparison $m_\mu/m_e$ vs.\ $\sqrt{43200}=120\sqrt{3}$ is \emph{not a
clean rational comparison}: $\sqrt{43200}$ is irrational, and the $-0.52\,\%$
residual is not an independent physical finding. It equals exactly half the
quadratic residual ($-1.034\,\%/2\approx-0.517\,\%$), because taking the square root
approximately halves the deviation. Nothing new becomes visible.

Only two comparisons are algebraically clean:
\begin{itemize}\itemsep2pt
\item \textbf{Square:} $(m_\mu/m_e)^2 = 43200$ --- both sides rational,
  residual $-1.034\,\%$ physical.
\item \textbf{Product:} $m_e\cdot m_\mu = 54~\text{MeV}^2$ --- both sides rational,
  residual $-0.016\,\%$ \textbf{[K]}{} (Doc.~338; algebraic origin:
  $|GF(3)^*|\cdot|GF(27)| = 2\cdot27 = 54$).
\end{itemize}

The product $m_e\cdot m_\mu = 54~\text{MeV}^2$ is the algebraically cleanest
Galois identity with the smallest rational residual ($-0.016\,\%$). The absolute
deviation is $-0.0087~\text{MeV}^2$ ($-$7386\,\sigma$ (CODATA) / $1386\,\sigma$ (Eides)$ CODATA, $-1386\,\sigma$ Eides) --- also physical, but
65 times smaller than the quadratic residual.

\textbf{Bookkeeping rule:} Galois identities should be compared only
in their natural algebraic form --- where no irrational numbers appear.
For $(m_\mu/m_e)^2 = 43200$ that is the square comparison; for $m_e\cdot m_\mu$
the product. The linear comparison $m_\mu/m_e$ vs.\ $\sqrt{43200}$ should be
avoided because it introduces an irrational intermediate
(verification script: \texttt{pruef\_351\_rationale\_vergleiche.py}).


\subsection*{Is the square comparison algebraically admissible?}

The identity $(m_\mu/m_e)^2=43200$ is admissible as a Galois comparison ---
both sides rational, no irrational factor. But it compares two different
objects:

\begin{itemize}\itemsep2pt
\item \textbf{Layer-1 (purely algebraic):} The identity
  $(r_\mu/r_e)^2/\xi=43200$ is exact from winding numbers ($r_\mu/r_e=12/5$,
  Doc.~317) and Galois group orders ($|GF(9)^*|^2\cdot5^2\cdot|GF(27)|$,
  Doc.~338). No measured value enters.
\item \textbf{Layer-2 (empirical):} $(m_\mu/m_e)^2_\text{PDG}=42753.12$
  comes from the HFS measurement (LAMPF 1999), extracted via $\nu_F\propto\alpha^2$.
  It is a QED-dependent measured value.
\end{itemize}

Following Stefaan Vossen's precision (5~Sept.~2026): the Layer-1 object
$(r_\mu/r_e)^2/\xi$ and the Layer-2 object $(m_\mu/m_e)^2_\text{PDG}$ are
\emph{distinguishable scientific objects}. The comparison is admissible, but
the residual $\approx78\cdot\xi$ belongs to the Layer-1~$\to$~Layer-2 mapping,
not to the algebraic identity itself. It lies far outside the SI projection
range ($\approx\xi$) and is open \SM.

The product $m_e\cdot m_\mu=54~\text{MeV}^2$ is the cleaner object in this
sense: its residual $\approx1.2\cdot\xi$ lies entirely within the range of the
SI projection approximation (Docs.~085, 306, 333) and is structurally
explainable \KM. The Layer-1~$\to$~Layer-2 transition is there covered by
the projection.






\subsection*{Two orthogonal axes: difference and sum of logarithms}

The PDG masses $m_e$ and $m_\mu$ produce two independent logarithmic projections
when compared with the Galois identities:

\[
\underbrace{\ln\!\frac{m_\mu/m_e}{\sqrt{43200}}}_{\text{difference axis}}
= \varepsilon_\mu - \varepsilon_e \approx -38.99\cdot\xi
\qquad
\underbrace{\ln\!\frac{m_e\cdot m_\mu}{54\,\text{MeV}^2}}_{\text{sum axis}}
= \varepsilon_e + \varepsilon_\mu \approx -1.21\cdot\xi
\]

\textbf{Difference axis: linear and square are one single finding.}
The linear ratio $m_\mu/m_e$ vs.\ $\sqrt{43200}$ ($-0.519\%$) and the
square $(m_\mu/m_e)^2$ vs.\ $43200$ ($-1.034\%$) lie on \emph{exactly the
same axis} $\varepsilon_\mu-\varepsilon_e$. Logarithmically it holds exactly:
\[
\ln\!\frac{(m_\mu/m_e)^2}{43200} = 2\cdot\ln\!\frac{m_\mu/m_e}{\sqrt{43200}}.
\]
The square doubles the percentage but measures no new physics.
It is \emph{one} finding, two ways of writing it.
The value $-38.99\cdot\xi$ is empirical; it lies $8\sigma$ (Eides) from
$39\cdot\xi$ and is not a Galois integer \SM.

\textbf{Sum axis: the product is orthogonal.}
$\varepsilon_e\approx+18.9\cdot\xi$ and $\varepsilon_\mu\approx-20.1\cdot\xi$
are nearly equal in magnitude and opposite in sign. In the sum they nearly
cancel; what remains is $-1.21\cdot\xi$ --- the common scale deviation of both
masses. This is at order $\xi$, consistent with the SI projection
$S^1_m\to\mathbb{R}_t$ (Docs.~085, 306, 333) \KM. The sum ($-1.21\cdot\xi$)
does not follow from the difference ($-38.99\cdot\xi$) --- they are two
independent physical pieces of information.

\begin{center}
\begin{tabular}{@{}llll@{}}
\toprule
Combination & Logarithm & Residual & Status \\
\midrule
$\ln(m_\mu/m_e\,/\sqrt{43200})$ & $\varepsilon_\mu-\varepsilon_e$ & $-38.99\cdot\xi$ & ratio residual, order $100\xi$ \KM \\
$\ln((m_\mu/m_e)^2/43200)$ & $2(\varepsilon_\mu-\varepsilon_e)$ & $-77.99\cdot\xi$ & exactly $=2\times$ above \\
$\ln(m_e\cdot m_\mu\,/\,54\,\text{MeV}^2)$ & $\varepsilon_e+\varepsilon_\mu$ & $-1.21\cdot\xi$ & derivation from decompactification \SM \\
\bottomrule
\end{tabular}
\end{center}

The apparently different percentages $-0.016\%$, $-0.52\%$, and $-1.03\%$
arise from two numbers: the generation tension ($-39\cdot\xi$) and the SI
projection ($-1.2\cdot\xi$). Their large ratio is not an amplification but
a cancellation in the sum
(verification script: \texttt{pruef\_351\_generationsdifferenz.py}).

The residual of the product comparison lies at order $\xi$ itself:

\begin{center}
\begin{tabular}{@{}lll@{}}
\toprule
Quantity & Residual in units of $\xi$ & Interpretation \\
\midrule
$m_e\cdot m_\mu$ vs.\ $54$ & $-1.21\cdot\xi$ & order $\xi$: SI projection \KM \\
$(m_\mu/m_e)^2$ vs.\ $43200$ & $-77.6\cdot\xi$ & higher order: fractal-recursive \KM \\
\bottomrule
\end{tabular}
\end{center}

The SI projection $S^1_m\to\mathbb{R}_t$ (Docs.~085, 306, 333) structurally
produces errors of order $\xi$. The product residual $-0.016\,\%\approx1.2\cdot\xi$
lies precisely in this range and is therefore structurally explainable as an
SI transition approximation \KM.
The quadratic residual $-1.034\,\%\approx78\cdot\xi$ lies in higher order
and is explained by the fractal-recursive correction (Docs.~295, 146, 338).

\textbf{Result:} The Galois identity $m_e\cdot m_\mu=54~\text{MeV}^2$ is
the only currently known identity whose residual lies at order $\xi$ ---
the structurally expected SI projection error. The residual is therefore not
unexplained; it corresponds to the approximation of the SI transition.


\subsection*{Constitutional description (after Stefaan Vossen, 5~Sept.~2026)}

The following four-layer formulation summarises the epistemic status of the FFGFT comparisons:

\begin{center}
\begin{tabular}{@{}p{3.5cm}p{9.5cm}@{}}
\toprule
Layer & Status \\
\midrule
\textbf{Layer-1} & Exact within rational arithmetic: $(m_\mu/m_e)^2=43200$ and $m_e\cdot m_\mu=54~\text{MeV}^2$ are integer Galois identities without approximation \BM. \\[4pt]
\textbf{Layer-2} & One empirical dimensional anchor: $m_\mu=105.66~\text{MeV}$ (PDG). All absolute predictions follow from it \KM. \\[4pt]
\textbf{Residual} & Empirical discrepancy with a structural explanation of approximately the correct order: $(m_\mu/m_e)^2$ deviates by $\approx78\cdot\xi$ from the PDG value (fractal-recursive correction, Doc.~295); $m_e\cdot m_\mu$ deviates by $\approx1.2\cdot\xi$ (SI projection approximation, Docs.~085, 306, 333) \KM. \\[4pt]
\textbf{Quantitative bridge} & The forward derivation of the exact residual from the $S^1_m\to\mathbb{R}_t$ non-isomorphism is not in the corpus. \textbf{Open [S].} \\
\bottomrule
\end{tabular}
\end{center}

The open [S] is not a gap in the pejorative sense but a precisely constituted limit: the theory states explicitly where its formal construction stops.


\subsection*{Methodological note (after Stefaan Vossen, 5~Sept.~2026)}

\begin{quote}
\emph{``Constitution does not tell the Galois identity what it is.
It tells us what a particular scientific use of that identity is presently entitled to mean.''}
\end{quote}

Two distinctions to be preserved in the record:

\begin{enumerate}\itemsep4pt
\item \textbf{$54$ as a finite-field integer $\neq$ $54~\text{MeV}^2$ as a dimensional representation.}
  The identity $|GF(3)^*|\cdot|GF(27)|=2\cdot27=54$ belongs to the dimensionless Layer-1 algebra \BM.
  $m_e\cdot m_\mu=54~\text{MeV}^2$ is a dimensional representation produced by the mapping into
  SI units (Layer-2). Both are distinguishable scientific objects;
  the unit MeV$^2$ belongs to the dimensional mapping, not to the dimensionless algebra alone.

\item \textbf{Numerical correspondence vs.\ physical attribution.}
  The residual $m_e\cdot m_\mu\approx54~\text{MeV}^2$ with deviation $\approx1.2\cdot\xi$
  is an established numerical correspondence \KM.
  The physical attribution of that residual to the SI projection
  $S^1_m\to\mathbb{R}_t$ (Docs.~085, 306, 333) remains bounded by the still-open
  quantitative bridge \SM.
  The correspondence stands; the attribution is a proposal, not a derivation.
\end{enumerate}


\subsection*{Methodological remark: relation vs.\ prediction}

The three percentages ($-0.519\%$, $-1.034\%$, $-0.016\%$) are
\emph{not tolerance values}. Tolerance arises only when a
\textbf{calculated} mass deviates from a \textbf{measured} one ---
which requires an independent FFGFT prediction for $m_e$ or $m_\mu$
individually.

The Galois identities $43200$ and $54~\text{MeV}^2$ are
\textbf{relations} between both masses, not individual predictions.
As long as no independent FFGFT value for $m_\mu$ (or $m_e$)
exists, none of the three percentages has a defined limit.

Moreover: every direct mass carries the fractal correction
$K_\text{frak}$ and the SI projection $S^1_m\to\mathbb{R}_t$.
In the ratio, $K_\text{frak}$ cancels; the SI corrections
\emph{subtract} (difference axis, $\approx -39\cdot\xi$).
In the product, $K^2$ remains; the SI corrections \emph{add}
(sum axis, $\approx -1.2\cdot\xi$). The different percentages
arise from this --- not from measurement errors.

No independent single prediction for $m_\mu$ without an input anchor
exists in the corpus. With $m_e$ as anchor (SETZUNG) there are two routes:
\begin{itemize}\itemsep2pt
\item \textbf{Square route:} $m_\mu = m_e\cdot\sqrt{43200} = m_e\cdot120\sqrt{3}$ ---
  irrational, deviation $+0.52\%$. Not a clean Layer-1 value.
\item \textbf{Product route:} $m_\mu = 54\,\text{MeV}^2/m_e$ --- rational,
  deviation $+0.016\%\approx1.2\cdot\xi$. The only rational tolerance
  value for $m_\mu$ at anchor $m_e$.
\end{itemize}
But: $m_{e,\text{PDG}}$ already carries $K_\text{frak}$ and the SI projection.
Hence $m_\mu^\text{pred}=54/m_{e,\text{PDG}}\approx m_{\mu,\text{bare}}(1-\varepsilon_e)$
and $m_{\mu,\text{PDG}}\approx m_{\mu,\text{bare}}(1+\varepsilon_\mu)$, so:
\[
\frac{m_\mu^\text{pred}}{m_{\mu,\text{PDG}}}-1
\approx -(\varepsilon_e+\varepsilon_\mu) \approx -1.2\cdot\xi .
\]
This is the sum axis --- the same relation as the product residual.
The fractal correction does not cancel but is contained in
$\varepsilon_e$ and $\varepsilon_\mu$ and adds in the product.
A true Layer-1 tolerance value for $m_\mu$ without any correction
does not exist, because every direct mass carries $K_\text{frak}$
and the SI projection.

\section{Error accumulation along the derivation chain}

Each step carries a small residual (recursion $\sim100\xi$, projection $\sim\xi$, extraction theory $10^{-7}$--$10^{-10}$). Already at $m_\mu$, the first step after the anchor $m_e$:

\begin{center}
\begin{tabular}{@{}ll@{}}
\toprule
Route $m_e\to m_\mu$ & Residual \\
\midrule
Galois bare: $m_\mu=\sqrt{43200}\,m_e$ & $+0.52\,\%$ \\
with $\sqrt{K_\text{frak}}$: $m_\mu=\sqrt{43200\cdot74/75}\,m_e$ & $-0.15\,\%$ \\
via $\alpha$: $E_0=\sqrt{\alpha/\xi}$, $m_\mu=E_0^2/m_e$ & $+1.37\,\%$ \\
\bottomrule
\end{tabular}
\end{center}

At $m_\tau$ (second step), $v$ (third, via $\xi^{3/2}$), each derivation accumulates further. To decide where a residual sits, convergence of empirical values onto the algebraic ones would have to be established --- it is not: $m_\mu/m_e$ (1999), $\alpha$ ($5\sigma$ tension), $\Delta m^2_\text{atm}$ ($2.4\sigma$ shift in 2024), $v$ (not measured).

\textbf{Bookkeeping:} For every reported prediction, state: (1)~the anchor, (2)~the number of derivation steps, (3)~corrections applied and where, (4)~the comparison value with its extraction method. ``FFGFT predicts $m_\mu$ to 0.5\,\%'' is incomplete; complete: ``From $m_e$ (PDG), in one step via the Galois identity, without $K_\text{frak}$, $m_\mu$ follows with $+0.52\,\%$ against the 1999 muonium HFS value.''

\section{Consequences for the corpus}

\textbf{No numerical value in the corpus changes.} What changes:
\begin{enumerate}\itemsep2pt
\item \textbf{$v$ is neither anchor nor observable.} Docs.~006, 319, R104: ``$v=246$~GeV (measured)'' $\to$ ``$v=246$~GeV (SM definition from $G_F$)''. Register R112.
\item \textbf{$m_\mu/m_e$ is QED-dependent.} ``model-independent'' $\to$ ``least model-dependent available value''; $\nu_F\propto\alpha^2$. Register R113.
\item \textbf{$\alpha$ has no unique measured value below 1~ppb.} Cite the comparison value as CODATA average with Rb/Cs tension.
\item \textbf{$m_\tau$:} two methods, PDG average 1776.93(9), FFGFT $0.4\sigma$ --- open.
\item \textbf{New open item \SM:} $\tau_\mu$ directly from FFGFT, without detour via $G_F$ and $v$ --- the only convention-free test of the weak sector.
\end{enumerate}

\section*{Sources}

All sources verified on 5~September 2026 against INSPIRE, APS, Nature, arXiv, or PubMed. Numerical values computed in exact rational arithmetic against CODATA 2022.

\begin{thebibliography}{99}\small
\bibitem{eberhart} M.~Eberhart, M.~Fael, A.~Keshavarzi, D.~Nomura, K.~Schönwald, M.~Steinhauser, T.~Teubner, J.~Wright: \emph{Muon lifetime and Fermi constant: an update}, arXiv:2607.02657 (2026).
\bibitem{mulan} V.~Tishchenko et al.\ (MuLan): Phys.\ Rev.\ D \textbf{87}, 052003 (2013).
\bibitem{sirlin} A.~Sirlin: Phys.\ Rev.\ D \textbf{22}, 971 (1980).
\bibitem{eides26} M.~I.~Eides: \emph{Muonium Hyperfine Splitting Uncertainty Revisited}, arXiv:2510.07281v2 (2026).
\bibitem{liu99} W.~Liu et al.: Phys.\ Rev.\ Lett.\ \textbf{82}, 711 (1999).
\bibitem{eides01} M.~I.~Eides, H.~Grotch, V.~A.~Shelyuto: Phys.\ Rep.\ \textbf{342}, 63 (2001).
\bibitem{morel} L.~Morel, Z.~Yao, P.~Cladé, S.~Guellati-Khélifa: Nature \textbf{588}, 61 (2020).
\bibitem{parker} R.~H.~Parker, C.~Yu, W.~Zhong, B.~Estey, H.~Müller: Science \textbf{360}, 191 (2018).
\bibitem{fan} X.~Fan, T.~G.~Myers, B.~A.~D.~Sukra, G.~Gabrielse: Phys.\ Rev.\ Lett.\ \textbf{130}, 071801 (2023).
\bibitem{aoyama} T.~Aoyama, T.~Kinoshita, M.~Nio: Atoms \textbf{7}, 28 (2019).
\bibitem{besiii} M.~Ablikim et al.\ (BESIII): Phys.\ Rev.\ D \textbf{90}, 012001 (2014).
\bibitem{belle2} I.~Adachi et al.\ (Belle II): Phys.\ Rev.\ D \textbf{108}, 032006 (2023).
\bibitem{kedr} V.~V.~Anashin et al.\ (KEDR): JETP Lett.\ \textbf{85}, 347 (2007).
\bibitem{pdg} PDG 2024, pdgLive S035M: $\tau$ mass average $1776.93\pm0.09$~MeV.
\bibitem{sturm} S.~Sturm, F.~Köhler, J.~Zatorski et al.: Nature \textbf{506}, 467 (2014).
\bibitem{koehler} F.~Köhler, S.~Sturm, A.~Kracke et al.: J.\ Phys.\ B \textbf{48}, 144032 (2015).
\bibitem{heisse} F.~Heiße et al.: Phys.\ Rev.\ Lett.\ \textbf{131}, 253002 (2023).
\bibitem{li} Q.~Li, C.~Xue, J.-P.~Liu et al.: Nature \textbf{560}, 582 (2018).
\bibitem{rosi} G.~Rosi, F.~Sorrentino, L.~Cacciapuoti, M.~Prevedelli, G.~M.~Tino: Nature \textbf{510}, 518 (2014).
\bibitem{nufit} I.~Esteban, M.~C.~Gonzalez-Garcia, M.~Maltoni, I.~Martinez-Soler, J.~P.~Pinheiro, T.~Schwetz: JHEP \textbf{12} (2024) 216.
\bibitem{codata} P.~J.~Mohr, D.~B.~Newell, B.~N.~Taylor, E.~Tiesinga: Rev.\ Mod.\ Phys.\ \textbf{97}, 025002 (2025).
\end{thebibliography}

\appendix
\section{Affected derivation chains}
\label{sec:chains}

Which FFGFT chain depends on which comparison value, and what each finding changes.

\begin{figure}[h]
\centering
\includegraphics[width=0.92\textwidth]{../figures/ffgft_ableitungsketten.pdf}
\caption*{Derivation chains: from $\xi$ to comparison values. Dashed boxes = external comparison values; dashed arrows = Galois identities.}
\end{figure}

\begin{center}
\begin{tabular}{@{}p{3cm}p{3.5cm}p{3cm}p{4cm}@{}}
\toprule
Finding & Affected documents & Numerical correction & Wording correction \\
\midrule
$v$ is SM definition (R112) & Docs.~006, 319, R104 & none & ``$v=246$~GeV (measured)'' $\to$ ``(SM definition from $G_F$)'' \\
$G_F$: $\Delta q$, $192\pi^3$, $\sqrt2$ & Docs.~006, R104, 319 & none & state $\Delta q=-0.44\,\%$; clarify with/without $\Delta q$ \\
$m_\mu/m_e$ QED-dependent (R113) & Docs.~338, 317, 011 & none & ``model-independent'' $\to$ ``least model-dependent'' \\
$\alpha$ Rb/Cs $5\sigma$ & Docs.~011, 070, 338 & none & cite CODATA average with tension \\
$m_\tau$ two methods & Docs.~317, 006 & none & cite BESIII + Belle II; PDG average 1776.93(9) \\
Neutrinos $\Delta m^2$ & Docs.~339, 340 & none & freeze NuFIT configuration \\
$G$ & Doc.~A145 & none & --- \\
\midrule
\textbf{Total} & & \textbf{none} & wording clarifications \\
\bottomrule
\end{tabular}
\end{center}

The residual $(m_\mu/m_e)^2 = 43200$ vs.\ 42753.12 (difference 446.9) is
$240\,000\,\sigma$ (CODATA) or $45\,000\,\sigma$ (real uncertainty, Eides) of the measurement uncertainty --- physical (R113).
No numerical value in the corpus changes.
